% Owner: Member 2 (CP). Uses the shared notation of docs/Problem.md and report/sections/ilpFormulation.tex.
\section{Constraint Programming Formulation}\label{sec:cp}

\subsection{Variables and Constraints}\label{sec:cp-model}
Let $G=(V,E)$ be the undirected input graph, $n=|V|$, $m=|E|$, and
$f:V\rightarrow\{1,\ldots,n\}$ a vertex-to-position labeling.
The bandwidth is $B(f)=\max_{\{u,v\}\in E}|f(u)-f(v)|$, with $B(f)=0$
when $E=\varnothing$. Let $UB_0=B(f_0)$ for $f_0(v)=v$.
We use integer position variables $p_v\in\{1,\ldots,n\}$, distance
variables $d_{uv}\in\{0,\ldots,n-1\}$, and a bandwidth variable
$B\in\{0,\ldots,UB_0\}$. The model is:
\begin{align}
\min\quad & B, \label{eq:cp-objective}\\
\text{s.t.}\quad & \operatorname{AllDifferent}(p_v:v\in V),\label{eq:cp-alldiff}\\
&d_{uv}=|p_u-p_v| &&\forall\{u,v\}\in E,\label{eq:cp-distance}\\
&B=\max_{\{u,v\}\in E}d_{uv} && E\ne\varnothing,\label{eq:cp-max}\\
&B=0 && E=\varnothing.\label{eq:cp-empty}
\end{align}
The source model has $n+m+1$ integer variables and one global
\texttt{AllDifferent} constraint, $m$ absolute-distance equalities,
and one bandwidth equality. These are \emph{source-model} counts;
each backend may introduce additional constraints during transformation.

\subsection{Correctness}\label{sec:cp-correctness}
Constraint~\eqref{eq:cp-alldiff} assigns distinct values from a
size-$n$ domain to $n$ vertices, so the $p_v$ form a permutation of
$\{1,\ldots,n\}$, and every permutation satisfies this constraint.
Constraints~\eqref{eq:cp-distance}--\eqref{eq:cp-empty} force $B$
to equal the bandwidth of that permutation, including the edgeless
case. Conversely, every valid labeling defines a feasible assignment
with the same objective. Since $UB_0$ is attained by the identity
labeling, the bound $B\le UB_0$ cannot eliminate an optimal solution.
Hence minimizing $B$ in this model is equivalent to the graph BMP.

\subsection{Implementations and Solution Validation}\label{sec:cp-backends}
\texttt{scripts/cp/cpSolver.py} implements the same formulation for
\texttt{cp\_ortools} (OR-Tools CP-SAT) and \texttt{cp\_cpo} (IBM CPLEX
CP Optimizer); these are two \emph{backends}, not two CP formulations.
OR-Tools uses \texttt{AddAllDifferent}, \texttt{AddAbsEquality}, and
\texttt{AddMaxEquality}. IBM CP Optimizer uses \texttt{all\_diff},
\texttt{abs}, and \texttt{max}. The identity assignment is provided as
a hint/starting point and is also retained as the checked initial
incumbent. Each run uses one worker, a declared random seed, and the
time remaining after model construction.

On a new incumbent the code extracts $[p_1,\ldots,p_n]$; the shared
checker independently verifies the permutation and recomputes the
bandwidth. A failed check is an error, not a feasible result.
\texttt{OPTIMAL} requires native proof of optimality and a matching
certified objective bound; an initialized run that reaches its deadline
without proof is \texttt{FEASIBLE}. The model does not infer UNSAT
from interruption. As with the ILP formulation, a runtime comparison
must specify which backend was used.

\subsection{Verification Scope and Limitations}\label{sec:cp-validation}
The repository contains CP workflow tests on small graphs (including
paths, cycles, complete, disconnected, and edgeless graphs), an
independent brute-force oracle, and JSON output checks.
\texttt{scripts/cp/checkSolvers.py} checks whether each backend solves
a four-vertex instance with independently known optimum. Passing the
small probe establishes only limited functional availability; it does
\emph{not} demonstrate that a commercial license permits the full
benchmark campaign. Record the large-instance licensing check and
actual solver versions in the campaign's \texttt{Environment.json}.
